%% ── Commitment & Merkle ──────────────────────────────────────────────────────
\newcommand{\para}{\textsf{pp}}
\newcommand{\cm}{\textsf{cm}}
\newcommand{\rt}{\textsf{rt}}
\newcommand{\cmmt}{\textsf{cm}_{\textsf{mt}}}
\newcommand{\merkleC}{\textsf{MT.Com}}
\newcommand{\merkleV}{\textsf{MT.Verify}}
\newcommand{\merkleO}{\textsf{MT.Open}}
\newcommand{\pedersen}{\textsf{P.CM}}
\newcommand{\ck}{\textsf{ck}}
\newcommand{\pedC}{\textsf{Ped.Com}}
\newcommand{\pedO}{\textsf{Ped.Open}}
\newcommand{\pedV}{\textsf{Ped.Verify}}

%% ── Algebra ──────────────────────────────────────────────────────────────────
\newcommand{\mat}{\mathbf}          % matrix: \mat{A}
\newcommand{\vect}{\mathbf}         % vector: \vect{x}
%% \enc is already defined by cryptocode's 'primitives' module (as encryption).
%% We redefine it here to mean the linear-code encoding operator instead.
\renewcommand{\enc}{\mathsf{Enc}}   % encoding operator (overrides cryptocode)

%% ── SNARK ────────────────────────────────────────────────────────────────────
\newcommand{\snarkP}{\textsf{SNARK.P}}
\newcommand{\snarkV}{\textsf{SNARK.V}}
\newcommand{\bPi}{{\bf \Pi}}
\newcommand{\setup}{{\sf Setup}}
\newcommand{\com}{{\sf Com}}
\newcommand{\prove}{{\sf Prove}}
\renewcommand{\verify}{{\sf Verify}}
\newcommand{\relation}{{\mathcal{R}}}
\newcommand{\cp}{{\mathsf{cp}}}
\renewcommand{\pk}{{\textsf{pk}}}
\renewcommand{\vk}{{\textsf{vk}}}
\renewcommand{\ck}{{\textsf{ck}}}



%% ── SNARK statement / witness (replaces \mathbb{x}, \mathbb{w}) ─────────────
%% \mathbb works only for uppercase.
%% We use \mathsf to distinguish SNARK instances from vector \vect{x} (\mathbf{x}).
\newcommand{\stmt}{\mathsf{x}}      % SNARK public statement
\newcommand{\wit}{\mathsf{w}}       % SNARK witness
\newcommand{\comwit}{\mathsf{u}}

%% ── Misc ─────────────────────────────────────────────────────────────────────
\newcommand{\fold}{\textsf{Fold}}
\newcommand{\Hash}{\textsf{Hash}}
\newcommand{\vecBindV}{\textsf{Vec.Bind}}
\newcommand{\cpLinkP}{\textsf{CPLink.Prove}}
\newcommand{\cpLinkV}{\textsf{CPLink.Verify}}
\newcommand{\protocol}{\textsf{LAMP}}   % protocol name
\newcommand{\fmMerkle}{\protocol_{\mathsf{Merkle}}}
\newcommand{\fmKZG}{\protocol_{\mathsf{KZG}}}
\newcommand{\fmCPLink}{\protocol_{\mathsf{CPLink}}}
\newcommand{\fmQA}{\protocol_{\mathsf{QA}}}

%% ── Sets of objects (avoid \mathcal{V} collision with verifier) ──────────────
%% Vector set is written explicitly in text to avoid notation collision.

%% ── Theorem environments ─────────────────────────────────────────────────────
\theoremstyle{plain}
\newtheorem{theorem}{Theorem}[section]
\newtheorem{lemma}[theorem]{Lemma}
\newtheorem{claim}[theorem]{Claim}
\newtheorem{corollary}[theorem]{Corollary}

\theoremstyle{definition}
\newtheorem{definition}[theorem]{Definition}

\theoremstyle{remark}
\newtheorem{remark}{Remark}

%% Some conference templates provide \Description for accessibility metadata.
\providecommand{\Description}[1]{}
